\chapter{2008年重庆卷（理）}

\section{单选题}

\begin{problem}
复数 \(1 + \frac{2}{\i^3} =\)（\quad）

\choices
    {\(1 + 2\i\)}
    {\(1 - 2\i\)}
    {\(-1\)}
    {\(3\)}
\end{problem}

\begin{answer}
A
\end{answer}

\begin{solution}
由 \(\i^{3}=-\i\)，得 \(\frac{2}{\i^{3}}=\frac{2}{-\i}\cdot\frac{\i}{\i}=\frac{2\i}{-\i^{2}}=2\i\)，所以
\[
1+\frac{2}{\i^{3}}=1+2\i
\]
故选 A.
\end{solution}

\begin{problem}
设 \(m\)，\(n\) 是整数，则“\(m\)，\(n\) 均为偶数”是“\(m + n\) 是偶数”的（\quad）

\choices
    {充分而不必要条件}
    {必要而不充分条件}
    {充要条件}
    {既不充分也不必要条件}
\end{problem}

\begin{answer}
A
\end{answer}

\begin{solution}
若 \(m\)，\(n\) 均为偶数，则存在整数 \(p\)，\(q\)，使得 \(m=2p\)，\(n=2q\)，从而 \(m+n=2(p+q)\) 为偶数. 因此“\(m\)，\(n\) 均为偶数”是“\(m+n\) 是偶数”的充分条件.

但当 \(m=n=1\) 时，\(m+n=2\) 为偶数，而 \(m\)，\(n\) 不均为偶数，所以它不是必要条件. 故选 A.
\end{solution}

\begin{problem}
圆 \(O_{1}: x^{2} + y^{2} - 2x = 0\) 和圆 \(O_{2}: x^{2} + y^{2} - 4y = 0\) 的位置关系是（\quad）

\choices
    {相离}
    {相交}
    {外切}
    {内切}
\end{problem}

\begin{answer}
B
\end{answer}

\begin{solution}
圆 \(O_{1}\) 的圆心为 \(O_{1}(1,0)\)，半径为 \(1\)；圆 \(O_{2}\) 的圆心为 \(O_{2}(0,2)\)，半径为 \(2\). 两圆心之间的距离为
\[
|O_{1}O_{2}|=\sqrt{(1-0)^{2}+(0-2)^{2}}=\sqrt{5}
\]
由于 \(2-1<\sqrt{5}<2+1\)，两圆相交. 故选 B.
\end{solution}

\begin{problem}
已知函数 \( y = \sqrt{1 - x} + \sqrt{x + 3} \) 的最大值为 \(M\)，最小值为 \( m \)，则 \( \frac{m}{M} \) 的值为（\quad）

\choices
    {\( \frac{1}{4} \)}
    {\( \frac{1}{2} \)}
    {\( \frac{\sqrt{2}}{2} \)}
    {\( \frac{\sqrt{3}}{2} \)}
\end{problem}

\begin{answer}
C
\end{answer}

\begin{solution}
函数的定义域为 \([-3,1]\). 令
\(u=\sqrt{1-x}\)，\(v=\sqrt{x+3}\)
则 \(u\geqslant0\)，\(v\geqslant0\)，且 \(u^{2}+v^{2}=4\)，函数值为 \(u+v\). 由柯西不等式，
\[
u+v\leqslant\sqrt{2(u^{2}+v^{2})}=2\sqrt{2}
\]
当 \(u=v=\sqrt{2}\) 时等号成立，所以最大值 \(M=2\sqrt{2}\). 又因为
\[
(u+v)^{2}=u^{2}+v^{2}+2uv\geqslant4
\]
且在 \(u=0\) 或 \(v=0\) 时等号成立，所以最小值 \(m=2\). 因此
\[
\frac{m}{M}=\frac{2}{2\sqrt{2}}=\frac{\sqrt{2}}{2}
\]
故选 C.
\end{solution}

\begin{problem}
已知随机变量 \(\xi\) 服从正态分布 \(N(3, \sigma^2)\)，则 \(P(\xi < 3) =\)（\quad）

\choices
    {\(\frac{1}{5}\)}
    {\( \frac{1}{4} \)}
    {\( \frac{1}{3} \)}
    {\(\frac{1}{2}\)}
\end{problem}

\begin{answer}
D
\end{answer}

\begin{solution}
正态分布 \(N(3,\sigma^{2})\) 关于均值 \(3\) 对称，且正态分布为连续型分布，因此
\[
P(\xi<3)=P(\xi>3)=\frac{1}{2}
\]
故选 D.
\end{solution}

\begin{problem}
若定义在 \(\R\) 上的函数 \(f(x)\) 满足：对任意 \(x_{1}\)，\(x_{2} \in \R\) 有 \(f(x_{1} + x_{2}) = f(x_{1}) + f(x_{2}) + 1\)，则下列说法一定正确的是（\quad）

\choices
    {\(f(x)\) 为奇函数}
    {\(f(x)\) 为偶函数}
    {\(f(x) + 1\) 为奇函数}
    {\(f(x) + 1\) 为偶函数}
\end{problem}

\begin{answer}
C
\end{answer}

\begin{solution}
令 \(g(x)=f(x)+1\)，则已知条件化为
\[
g(x_{1}+x_{2})=f(x_{1}+x_{2})+1=f(x_{1})+f(x_{2})+2=g(x_{1})+g(x_{2})
\]
取 \(x_{1}=x_{2}=0\)，得 \(g(0)=0\). 再取 \(x_{1}=x\)，\(x_{2}=-x\)，得
\[
0=g(0)=g(x)+g(-x)
\]
所以 \(g(-x)=-g(x)\)，即 \(f(x)+1\) 为奇函数. 故选 C.
\end{solution}

\begin{problem}
若过两点 \(P_{1}(-1,2)\)，\(P_{2}(5,6)\) 的直线与 \(x\) 轴相交于点 \(P\)，则点 \(P\) 分有向线段 \(\overrightarrow{P_1P_2}\) 所成的比 \(\lambda\) 的值为（\quad）

\choices
    {\(-\frac{1}{3}\)}
    {\(-\frac{1}{5}\)}
    {\(\frac{1}{5}\)}
    {\(\frac{1}{3}\)}
\end{problem}

\begin{answer}
A
\end{answer}

\begin{solution}
直线 \(P_{1}P_{2}\) 的斜率为
\[
\frac{6-2}{5-(-1)}=\frac{2}{3}
\]
其方程为 \(y-2=\frac{2}{3}(x+1)\). 令 \(y=0\)，解得 \(x=-4\)，所以 \(P=(-4,0)\). 于是
\(\overrightarrow{P_{1}P}=(-3,-2)\)，\(\overrightarrow{PP_{2}}=(9,6)\)
从而 \(\overrightarrow{P_{1}P}=-\frac{1}{3}\overrightarrow{PP_{2}}\)，故 \(\lambda=-\frac{1}{3}\). 故选 A.
\end{solution}

\begin{problem}
已知双曲线 \(\frac{x^2}{a^2} - \frac{y^2}{b^2} = 1\)（\(a > 0\)，\(b > 0\)）的一条渐近线为 \(y = kx\)（\(k > 0\)），离心率 \(e = \sqrt{5} k\)，则双曲线方程为（\quad）

\choices
    {\(\frac{x^2}{a^2} - \frac{y^2}{4a^2} = 1\)}
    {\(\frac{x^2}{a^2} - \frac{y^2}{5a^2} = 1\)}
    {\(\frac{x^2}{4b^2} - \frac{y^2}{b^2} = 1\)}
    {\(\frac{x^2}{5b^2} - \frac{y^2}{b^2} = 1\)}
\end{problem}

\begin{answer}
C
\end{answer}

\begin{solution}
双曲线的渐近线斜率满足 \(k=\frac{b}{a}\)，离心率满足
\[
e=\frac{\sqrt{a^{2}+b^{2}}}{a}=\sqrt{1+k^{2}}
\]
由 \(e=\sqrt{5}k\) 且 \(k>0\)，得
\(1+k^{2}=5k^{2}\)，\(k=\frac{1}{2}\).
所以 \(a=2b\)，原双曲线方程可写为
\[
\frac{x^{2}}{4b^{2}}-\frac{y^{2}}{b^{2}}=1
\]
故选 C.
\end{solution}

\begin{problem}
如图，体积为 \(V\) 的大球内有 \(4\) 个小球，每个小球的球面过大球球心且与大球球面有且只有一个交点，\(4\) 个小球的球心是以大球球心为中心的正方形的 \(4\) 个顶点. \(V_{1}\) 为小球相交部分 （图中阴影部分） 的体积，\(V_{2}\) 为大球内，小球外的图中黑色部分的体积，则下列关系中正确的是（\quad）

\bitmapfigure[float=inline,max width=0.42\linewidth]{img/2008/chongqing_science/q9_fig1.png}
\FloatBarrier

\choices
    {\(V_{1} > \frac{V}{2}\)}
    {\(V_{2} < \frac{V}{2}\)}
    {\(V_{1} > V_{2}\)}
    {\(V_{1} < V_{2}\)}
\end{problem}

\begin{answer}
D
\end{answer}

\begin{solution}
设大球半径为 \(R\)，小球半径为 \(r\)，小球球心为 \(S\)，大球球心为 \(O\). 因为小球面过 \(O\)，所以 \(|OS|=r\)；又因为小球内切于大球，且两球面只有一个交点，所以 \(|OS|=R-r\). 因此 \(r=\frac{R}{2}\).

每个小球的体积为

\[
\frac{4}{3}\pi r^{3}=\frac{4}{3}\pi\left(\frac{R}{2}\right)^{3}=\frac{V}{8}
\]

四个小球体积之和为 \(V/2\). 正方形的两个相对顶点所对应的小球球心关于 \(O\) 对称，故这两个球心之间的距离为 \(R=2r\)，相应的两个小球除 \(O\) 外不相交. 任意三个小球都包含一对相对小球，因此它们的公共部分至多为点 \(O\). 因而图中四块阴影分别是四对相邻小球的交叠部分，且这些交叠部分在体积意义下互不重叠；其总体积为 \(V_{1}\). 四个小球的并集体积为

\[
\frac{V}{2}-V_{1}
\]

所以大球内小球外部分的体积为

\[
V_{2}=V-\left(\frac{V}{2}-V_{1}\right)=\frac{V}{2}+V_{1}>V_{1}
\]

故选 D.
\end{solution}

\begin{problem}
函数 \( f(x) = \frac{\sin x - 1}{\sqrt{3 - 2\cos x - 2\sin x}} \) 的值域是（\quad）

\choices
    {\( \left[-\frac{\sqrt{2}}{2},0\right] \)}
    {\([-1,0]\)}
    {\([- \sqrt{2}, 0]\)}
    {\([- \sqrt{3}, 0]\)}
\end{problem}

\begin{answer}
B
\end{answer}

\begin{solution}
记 \(s=\sin x\)，\(c=\cos x\). 分母的平方满足
\[
3-2c-2s\geqslant3-2\sqrt{2}>0
\]
所以函数处处有定义，且分子 \(s-1\leqslant0\)，从而 \(f(x)\leqslant0\). 又有
\[
\bigl(3-2c-2s\bigr)-(1-s)^{2}
=2-2c-s^{2}
=(1-c)^{2}\geqslant0
\]
因此
\[
0\leqslant\frac{(1-s)^{2}}{3-2c-2s}\leqslant1
\]
即 \(-1\leqslant f(x)\leqslant0\). 当 \(x=0\) 时 \(f(x)=-1\)，当 \(x=\frac{\pi}{2}\) 时 \(f(x)=0\). 函数在 \([0,\frac{\pi}{2}]\) 上连续，所以能取得这两个端点之间的所有值，值域为 \([-1,0]\). 故选 B.
\end{solution}

\section{填空题}

\begin{problem}
设集合 \(U = \{1, 2, 3, 4, 5\}\)，\(A = \{2, 4\}\)，\(B = \{3, 4, 5\}\)，\(C = \{3, 4\}\)，则 \((A \cup B) \cap (\complement_U C) =\fillinblank{}\).
\end{problem}

\begin{answer}
\(\{2,5\}\)
\end{answer}

\begin{solution}
由 \(A\cup B=\{2,3,4,5\}\)，以及 \(\complement_U C=U\smallsetminus C=\{1,2,5\}\)，得
\[
(A\cup B)\cap(\complement_U C)=\{2,3,4,5\}\cap\{1,2,5\}=\{2,5\}
\]
\end{solution}

\begin{problem}
已知函数 \( f(x) = \begin{cases}
2x + 3 & (\text{当 } x \neq 0 \text{ 时})\\
a & (\text{当 } x = 0 \text{ 时})
\end{cases} \) 在点 \( x = 0 \) 处连续，则 \( \displaystyle\lim_{n \to \infty} \frac{a n^{2} + 1}{a^{2} n^{2} + n} = \)\fillinblank{}.
\end{problem}

\begin{answer}
\(\frac{1}{3}\)
\end{answer}

\begin{solution}
由函数在 \(x=0\) 处连续，得
\[
a=f(0)=\lim_{x\to0}(2x+3)=3
\]
因此
\[
\displaystyle\lim_{n\to\infty}\frac{an^{2}+1}{a^{2}n^{2}+n}
=\displaystyle\lim_{n\to\infty}\frac{a+n^{-2}}{a^{2}+n^{-1}}
=\frac{3}{9}=\frac{1}{3}
\]
\end{solution}

\begin{problem}
已知 \( a^{\frac{2}{3}} = \frac{4}{9}\)（\(a > 0\)），则 \( \log_{\frac{2}{3}} a = \)\fillinblank{}.
\end{problem}

\begin{answer}
3
\end{answer}

\begin{solution}
因为 \(a>0\)，由
\[
a^{\frac{2}{3}}=\left(\frac{2}{3}\right)^{2}
\]
可得 \(a^{1/3}=\frac{2}{3}\)，从而 \(a=\left(\frac{2}{3}\right)^{3}\). 所以
\[
\log_{\frac{2}{3}}a=3
\]
\end{solution}

\begin{problem}
设 \(S_{n}\) 是等差数列 \(\{a_{n}\}\) 的前 \(n\) 项和，\(a_{12} = -8\)，\(S_{9} = -9\). 则 \(S_{16} =\fillinblank{}\).
\end{problem}

\begin{answer}
\(-72\)
\end{answer}

\begin{solution}
设等差数列的首项为 \(a_{1}\)，公差为 \(d\). 由已知条件，
\[
a_{1}+11d=-8
\]
以及
\[
S_{9}=\frac{9}{2}(2a_{1}+8d)=-9
\]
所以 \(a_{1}+4d=-1\). 两式相减得 \(7d=-7\)，即 \(d=-1\)，进而 \(a_{1}=3\). 因此
\[
S_{16}=\frac{16}{2}\bigl(2a_{1}+15d\bigr)=8(6-15)=-72
\]
\end{solution}

\begin{problem}
直线 \(l\) 与圆 \(x^{2} + y^{2} + 2x - 4y + a = 0\)（\(a < 3\)）相交于两点 \(A\)，\(B\)，弦 \(AB\) 的中点为 \((0, 1)\)，则直线 \(l\) 的方程为\fillinblank{}.
\end{problem}

\begin{answer}
\(y=x+1\)
\end{answer}

\begin{solution}
圆的圆心为 \(O(-1,2)\)，半径的平方为 \(5-a\). 圆心与弦中点 \(Q(0,1)\) 的连线斜率为
\[
\frac{1-2}{0-(-1)}=-1
\]
圆心到弦中点的半径垂直于弦，所以直线 \(l\) 的斜率为 \(1\). 又直线经过 \((0,1)\)，故
\[
l:y=x+1
\]
由于 \(a<3\)，圆的半径大于 \(\sqrt{2}\)，而圆心到该直线的距离为 \(\sqrt{2}\)，所以它确实与圆相交于两点.
\end{solution}

\begin{problem}
某人有 \(4\) 种颜色的灯泡（每种颜色的灯泡足够多），要在如图所示的 \(6\) 个点 \(A\)，\(B\)，\(C\)， \(A_{1}\)， \(B_{1}\)， \(C_{1}\) 上各装一个灯泡，要求同一条线段两端的灯泡不同色. 则每种颜色的灯泡都至少用一个的安装方法共有\fillinblank{} 种.

\bitmapfigure[float=inline,max width=0.38\linewidth]{img/2008/chongqing_science/q16_fig1.png}
\FloatBarrier

\end{problem}

\begin{answer}
216
\end{answer}

\begin{solution}
把四种颜色分别记为四个不同的颜色. 三角形 \(ABC\) 的三个顶点两两相连，因此其颜色必须互不相同，颜色安排有
\[
4\cdot3\cdot2=24
\]
种. 固定 \(A\)，\(B\)，\(C\) 的颜色后，设对应的颜色为 \(1\)，\(2\)，\(3\)，第四种颜色为 \(4\). 三角形 \(A_{1}B_{1}C_{1}\) 的三个顶点也必须两两异色，若暂不考虑三条竖直线段的限制，共有 \(4\cdot3\cdot2=24\) 种安排.

设事件 \(E_{A}\)，\(E_{B}\)，\(E_{C}\) 分别表示 \(A_{1}\) 与 \(A\)，\(B_{1}\) 与 \(B\)，\(C_{1}\) 与 \(C\) 同色. 则每个事件有 \(3\cdot2=6\) 种安排；任意两个事件同时发生时，第三个顶点有 \(2\) 种颜色可选；三个事件同时发生时有 \(1\) 种安排. 因此满足所有线段两端异色的底面安排数为
\[
24-3\cdot6+3\cdot2-1=11
\]
满足相邻顶点异色的总安排数为 \(24\cdot11=264\).

还要每种颜色至少使用一次. 若只使用三种颜色，先从四种颜色中选出未使用的一种，有 \(4\) 种选法；固定这三种颜色后，顶面三角形有 \(3!=6\) 种排列，底面三角形必须相对于顶面三个对应位置构成三阶错排，有 \(2\) 种排列. 因此只使用三种颜色的安排数为 \(4\cdot6\cdot2=48\). 所以所求安装方法数为
\[
264-48=216
\]
\end{solution}

\section{解答题}

\begin{problem}
设 \(\triangle ABC\) 的内角 \(A\)，\(B\)，\(C\) 的对边分别为 \(a\)，\(b\)，\(c\)，且 \(A = 60^{\circ}\)，\(c = 3b\). 求：
\begin{enumerate}
\item \(\frac{a}{c}\) 的值；
\item \(\cot B + \cot C\) 的值.
\end{enumerate}
\end{problem}

\begin{solution}
由余弦定理，
\[
a^{2}=b^{2}+c^{2}-2bc\cos A=b^{2}+9b^{2}-2b\cdot3b\cdot\frac{1}{2}=7b^{2}
\]
由于 \(c=3b\)，所以
\[
\frac{a}{c}=\frac{\sqrt{7}b}{3b}=\frac{\sqrt{7}}{3}
\]

设三角形面积为 \(S\). 由面积公式，
\[
S=\frac{1}{2}bc\sin A=\frac{1}{2}\cdot b\cdot3b\cdot\frac{\sqrt{3}}{2}=\frac{3\sqrt{3}}{4}b^{2}
\]
又由余弦定理的变形，
\(\cot B=\frac{a^{2}+c^{2}-b^{2}}{4S}\)，\(\cot C=\frac{a^{2}+b^{2}-c^{2}}{4S}\).
因此
\[
\cot B+\cot C=\frac{2a^{2}}{4S}=\frac{7b^{2}}{2S}=\frac{7b^{2}}{2\cdot\frac{3\sqrt{3}}{4}b^{2}}=\frac{14\sqrt{3}}{9}
\]
\end{solution}

\begin{problem}
甲，乙，丙三人按下面的规则进行乒乓球比赛：第一局由甲，乙参加而丙轮空，以后每一局由前一局的获胜者与轮空者进行比赛，而前一局的失败者轮空. 比赛按这种规则一直进行到其中一人连胜两局或打满 \(6\) 局时停止. 设在每局中参赛者胜负的概率均为 \(\frac{1}{2}\)，且各局胜负相互独立. 求：
\begin{enumerate}
\item 打满 \(3\) 局比赛还未停止的概率；
\item 比赛停止时已打局数 \( \xi \) 的分布列与期望 \(E(\xi)\).
\end{enumerate}
\end{problem}

\begin{solution}
第一局结束后，下一局由上一局的获胜者与轮空者比赛. 若上一局的获胜者再次获胜，则该人连胜两局，比赛停止；若轮空者获胜，则比赛继续. 由于每局两名参赛者胜负概率均为 \(\frac{1}{2}\)，每次比赛继续的概率也是 \(\frac{1}{2}\).

打满三局而比赛尚未停止，等价于第二局和第三局都由上一局获胜者的对手获胜，因此其概率为
\[
\left(\frac{1}{2}\right)^{2}=\frac{1}{4}
\]

比赛在第 \(2\)，\(3\)，\(4\)，\(5\) 局停止，分别需要在此前连续继续，并在该局由上一局获胜者再次获胜；比赛打到第 \(6\) 局则无论该局结果如何都停止. 因此
\(P(\xi=2)=\frac{1}{2}\)，\(P(\xi=3)=\frac{1}{4}\)，\(P(\xi=4)=\frac{1}{8}\)，\(P(\xi=5)=\frac{1}{16}\)，\(P(\xi=6)=\frac{1}{16}\).
分布列为
\begin{center}
\begin{tabular}{c|ccccc}
\(\xi\) & \(2\) & \(3\) & \(4\) & \(5\) & \(6\) \\
\hline
\(P\) & \(\frac{1}{2}\) & \(\frac{1}{4}\) & \(\frac{1}{8}\) & \(\frac{1}{16}\) & \(\frac{1}{16}\)
\end{tabular}
\end{center}
所以
\[
E(\xi)=2\cdot\frac{1}{2}+3\cdot\frac{1}{4}+4\cdot\frac{1}{8}+5\cdot\frac{1}{16}+6\cdot\frac{1}{16}=\frac{47}{16}
\]
\end{solution}

\begin{problem}
如图，在 \(\triangle ABC\) 中，\(B = 90^{\circ}\)，\(AC = \frac{15}{2}\)，\(D\)，\(E\) 两点分别在 \(AB\)，\(AC\) 上，使 \(\frac{AD}{DB} = \frac{AE}{EC} = 2\)，\(DE = 3\). 现将 \(\triangle ABC\) 沿 \(DE\) 折成直二面角，求：
\begin{enumerate}
\item 异面直线 \(AD\) 与 \(BC\) 的距离；
\item 二面角 \(A - EC - B\) 的大小（用反三角函数表示）.
\end{enumerate}

\FigureLayoutDeclare{img/2008/chongqing_science/q19_fig1.png}{12.0cm}{27bp 53bp 1bp 4bp}
\bitmapfigure[float=inline,max width=0.62\linewidth]{img/2008/chongqing_science/q19_fig1.png}
\FloatBarrier
\end{problem}

\begin{solution}
因为 \(D\)，\(E\) 分别将 \(AB\)，\(AC\) 按相同比例分割，所以 \(DE\parallel BC\)，且
\[
\frac{DE}{BC}=\frac{AD}{AB}=\frac{2}{3}
\]
由 \(DE=3\)，得 \(BC=\frac{9}{2}\). 在直角三角形 \(ABC\) 中，
\[
AB=\sqrt{AC^{2}-BC^{2}}=\sqrt{\left(\frac{15}{2}\right)^{2}-\left(\frac{9}{2}\right)^{2}}=6
\]
所以 \(AD=4\)，\(DB=2\)，\(AE=5\)，\(EC=\frac{5}{2}\).

折叠后建立空间直角坐标系，使折痕 \(DE\) 为 \(x\) 轴，取
\(D(0,0,0)\)，\(E(3,0,0)\)，\(A(0,4,0)\)，\(B(0,0,-2)\)，\(C\left(\frac{9}{2},0,-2\right)\).
这组坐标满足两侧平面分别为 \(z=0\) 和 \(y=0\)，且二面角为直二面角.

直线 \(AD\) 的方向向量可取 \((0,1,0)\)，直线 \(BC\) 的方向向量可取 \((1,0,0)\). 又 \(D\in AD\)，\(B\in BC\)，且
\[
\overrightarrow{DB}=(0,0,-2)
\]
同时垂直于这两条直线，所以异面直线 \(AD\) 与 \(BC\) 的距离为 \(DB=2\).

在平面 \(AEC\) 中，取
\(\overrightarrow{EA}=(-3,4,0)\)，\(\overrightarrow{EC}=\left(\frac{3}{2},0,-2\right)\)
其法向量可取
\[
\bs{n}_{1}=\overrightarrow{EA}\times\overrightarrow{EC}=(-8,-6,-6)
\]
平面 \(BEC\) 为 \(y=0\)，可取其法向量 \(\bs{n}_{2}=(0,1,0)\). 设二面角 \(A-EC-B\) 的平面角为 \(\theta\)，则
\[
\cos\theta=\frac{|\bs{n}_{1}\cdot\bs{n}_{2}|}{|\bs{n}_{1}|\,|\bs{n}_{2}|}=\frac{6}{2\sqrt{34}}=\frac{3}{\sqrt{34}}
\]
由于二面角取其平面角的锐角，
\[
\tan\theta=\sqrt{\frac{1}{\cos^{2}\theta}-1}=\sqrt{\frac{34}{9}-1}=\frac{5}{3}
\]
故二面角的大小为 \(\arctan\frac{5}{3}\).
\end{solution}

\begin{problem}
设函数 \( f(x) = ax^2 + bx + c\)（\(a \neq 0\)），曲线 \( y = f(x) \) 通过点 \((0, 2a + 3)\)，且在点 \((-1, f(-1))\) 处的切线垂直于 \(y\) 轴.
\begin{enumerate}
\item 用 \(a\) 分别表示 \(b\) 和 \(c\)；
\item 当 \(bc\) 取得最小值时，求函数 \(g(x) = -f(x)\e^{-x}\) 的单调区间.
\end{enumerate}
\end{problem}

\begin{solution}
由曲线经过 \((0,2a+3)\)，得
\[
c=f(0)=2a+3
\]
又因为切线垂直于 \(y\) 轴，所以该切线为水平切线，故
\[
f'(-1)=2a(-1)+b=0
\]
从而 \(b=2a\).

于是
\[
bc=2a(2a+3)=4a^{2}+6a=4\left(a+\frac{3}{4}\right)^{2}-\frac{9}{4}
\]
由于 \(a\neq0\)，且 \(a=-\frac{3}{4}\) 合法，所以 \(bc\) 的最小值在 \(a=-\frac{3}{4}\) 时取得. 此时
\(b=-\frac{3}{2}\)，\(c=\frac{3}{2}\)
从而
\(f(x)=-\frac{3}{4}x^{2}-\frac{3}{2}x+\frac{3}{2}\)，\(g(x)=\frac{3}{4}(x^{2}+2x-2)\e^{-x}\).
求导得
\[
g'(x)=\frac{3}{4}\bigl[(2x+2)-(x^{2}+2x-2)\bigr]\e^{-x}
=\frac{3}{4}(2-x)(x+2)\e^{-x}
\]
因为 \(\e^{-x}>0\)，所以当 \(x<-2\) 或 \(x>2\) 时 \(g'(x)<0\)，当 \(-2<x<2\) 时 \(g'(x)>0\). 因此 \(g(x)\) 在 \((-\infty,-2)\)，\((2,+\infty)\) 上单调递减，在 \((-2,2)\) 上单调递增.
\end{solution}

\begin{problem}
如图，\(M(-2,0)\) 和 \(N(2,0)\) 是平面上的两点，动点 \(P\) 满足： \(|PM| + |PN| = 6\).

\begin{enumerate}
\item 求点 \(P\) 的轨迹方程；
\item 若 \( |PM| \cdot |PN| = \frac{2}{1 - \cos MPN} \)，求点 \(P\) 的坐标.
\end{enumerate}

\bitmapfigure[float=inline,max width=0.42\linewidth]{img/2008/chongqing_science/q21_fig1.png}
\FloatBarrier
\end{problem}

\begin{solution}
由 \(|PM|+|PN|=6\)，可知点 \(P\) 的轨迹是以 \(M(-2,0)\)，\(N(2,0)\) 为焦点的椭圆. 设其长半轴为 \(a_{0}\)，焦距的一半为 \(c_{0}\)，则
\(2a_{0}=6\)，\(c_{0}=2\)，\(a_{0}=3\).
于是短半轴满足 \(b_{0}^{2}=a_{0}^{2}-c_{0}^{2}=9-4=5\)，故轨迹方程为
\[
\frac{x^{2}}{9}+\frac{y^{2}}{5}=1
\]

设 \(r=|PM|\)，\(s=|PN|\)，\(\theta=\angle MPN\)，并记 \(p=rs\). 由余弦定理及 \(|MN|=4\)，
\[
16=r^{2}+s^{2}-2rs\cos\theta=(r+s)^{2}-2p-2p\cos\theta
\]
因为 \(r+s=6\)，所以
\[
\cos\theta=\frac{r^{2}+s^{2}-16}{2p}=\frac{36-2p-16}{2p}=\frac{10}{p}-1
\]
题设条件给出 \(p(1-\cos\theta)=2\)，代入上式得
\(p\left(2-\frac{10}{p}\right)=2\)，\(p=6\).
因此 \(r\)，\(s\) 是方程
\[
t^{2}-6t+6=0
\]
的两个根，故 \(r+s=6\)，\(rs=6\)，并且
\[
(r-s)^{2}=(r+s)^{2}-4rs=36-24=12
\]
设 \(P=(x,y)\)，则
\[
r^{2}-s^{2}=\bigl((x+2)^{2}+y^{2}\bigr)-\bigl((x-2)^{2}+y^{2}\bigr)=8x
\]
另一方面，\(r^{2}-s^{2}=(r-s)(r+s)=\pm12\sqrt{3}\)，所以
\[
x=\pm\frac{3\sqrt{3}}{2}
\]
又
\[
r^{2}+s^{2}=(r+s)^{2}-2rs=24=2(x^{2}+y^{2}+4)
\]
从而 \(x^{2}+y^{2}=8\)，故 \(y^{2}=\frac{5}{4}\). 所以点 \(P\) 的坐标为
\(\left(\frac{3\sqrt{3}}{2},\frac{\sqrt{5}}{2}\right)\)，\(\left(\frac{3\sqrt{3}}{2},-\frac{\sqrt{5}}{2}\right)\)，\(\left(-\frac{3\sqrt{3}}{2},\frac{\sqrt{5}}{2}\right)\)，\(\left(-\frac{3\sqrt{3}}{2},-\frac{\sqrt{5}}{2}\right)\).
\end{solution}

\begin{problem}
设各项均为正数的数列 \(\{a_{n}\}\) 满足 \(a_1 = 2\)，\(a_n = a_{n+1}^{\frac{3}{2}} a_{n+2}\)（\(n \in \N^*\)）.
\begin{enumerate}
\item 若 \(a_2 = \frac{1}{4}\)，求 \(a_3\)，\(a_4\) 并猜想 \(a_{2008}\) 的值 （不需证明）；
\item 记 \( b_n = a_1a_2 \cdots a_n \)（\( n \in \N^* \)），若 \( b_n \geqslant 2\sqrt{2} \) 对 \( n \geqslant 2 \) 恒成立，求 \( a_2 \) 的值及数列 \( \{b_n\} \) 的通项公式.
\end{enumerate}
\end{problem}

\begin{solution}
令 \(x_n=\log_{2}a_n\). 由于各项均为正数，原递推式等价于
\[
x_n=\frac{3}{2}x_{n+1}+x_{n+2}
\]
即
\[
x_{n+2}=x_n-\frac{3}{2}x_{n+1}
\]

当 \(a_2=\frac14\) 时，\(x_1=1\)，\(x_2=-2\)，所以
\(x_3=x_1-\frac{3}{2}x_2=4\)，\(x_4=x_2-\frac{3}{2}x_3=-8\).
故 \(a_3=2^{4}=16\)，\(a_4=2^{-8}=\frac{1}{256}\). 此时数列 \(\{x_n\}\) 为 \(1\)，\(-2\)，\(4\)，\(-8\)，\(\ldots\). 其特征方程为
\[
\lambda^{2}+\frac{3}{2}\lambda-1=0
\]
根为 \(\frac12\)，\(-2\). 由初始值可得 \(x_n=(-2)^{n-1}\)，因此猜想
\[
a_{2008}=2^{x_{2008}}=2^{-2^{2007}}
\]

再设 \(t=x_2\). 递推数列的通项可写为
\[
x_n=u\left(\frac12\right)^{n-1}+v(-2)^{n-1}
\]
由 \(x_1=1\)，\(x_2=t\)，得
\(u+v=1\)，\(\frac{u}{2}-2v=t\)
从而
\(u=\frac{2(t+2)}{5}\)，\(v=\frac{1-2t}{5}\).
令 \(s_n=\log_{2}b_n=x_1+x_2+\cdots+x_n\)，则
\[
s_n=2u(1-2^{-n})+\frac{v}{3}\bigl(1-(-2)^n\bigr)
\]
若 \(v>0\)，则当 \(n\) 取偶数并趋于无穷时，\(s_n\) 趋于负无穷；若 \(v<0\)，则当 \(n\) 取奇数并趋于无穷时，\(s_n\) 趋于负无穷. 这两种情况都不可能满足对所有 \(n\geqslant2\) 有 \(b_n\geqslant2\sqrt{2}=2^{3/2}\)，故必须有 \(v=0\). 于是
\(t=\frac12\)，\(a_2=2^t=\sqrt{2}\).
此时 \(u=1\)，所以
\[
s_n=\sum_{k=1}^{n}2^{-(k-1)}=2(1-2^{-n})=2-2^{1-n}
\]
从而
\[
b_n=2^{s_n}=2^{2-2^{1-n}}
\]
当 \(n\geqslant2\) 时，\(s_n\geqslant s_2=\frac32\)，确实满足题设不等式.
\end{solution}
